% Paper 4, §3 — Lens techniques: the cheapest readout
% Thesis: lenses project intermediate hidden states through the model's
% own unembedding W_U to obtain a layer-l "as-if-final" distribution
% over the vocabulary. Logit lens (nostalgebraist 2020) is the no-training
% baseline; tuned lens (Belrose 2023) inserts a per-layer affine
% translator that fixes documented cross-model failures; DoLa (Chuang
% 2023) turns the framework into a hallucination-suppression decoding
% strategy.

\section{Lens techniques}
\label{sec:logit-lens-family}
\label{sec:lens-frontier}

A \emph{lens} takes a \glsterm{transformer}{transformer}'s intermediate residual-stream state
$\mathbf{h}_\ell$ at some layer $\ell$ and returns a distribution over
the \glsterm{vocabulary}{vocabulary} $\mathcal{V}$ — the prediction the model would commit to
if it were forced to early-exit at $\ell$. Lenses are the cheapest
interpretability readout: no training data, no auxiliary classifier, no
causal intervention; the lens reuses the model's own pre-unembed
normalization and unembedding matrix $W_U$. This section covers three
lens techniques that anchor the literature: the \emph{\glsterm{logit-lens}{logit lens}}
\citep{nostalgebraist2020-logit-lens}, the no-training baseline that
introduced the framework on a LessWrong post in 2020; the \emph{tuned
lens} \citep{belrose2023-tuned-lens}, which inserts a learned per-layer
affine translator and fixes the \glsterm{logit-lens}{logit lens}'s documented cross-model
failures; and \emph{\glsterm{dola}{DoLa}} \citep{chuang2023-dola}, an inference-time
contrastive-decoding strategy that uses the layerwise-readout structure
to suppress hallucinations. Each commits to a different question: what
\emph{would} the model predict at layer $\ell$? what \emph{is} the model
already predicting there once basis drift is accounted for? and what
content did the \emph{later} layers specifically commit?

\subsection{Setup: the iterative-inference view}
\label{sec:lens-setup}

\citet{belrose2023-tuned-lens} frame the \glsterm{transformer}{transformer} as iterative
inference\footnote{KB excerpt: \texttt{kb/excerpts/belrose2023-tuned-lens.md\#sec-1-iterative}.}:
each block makes a small additive update to the \glsterm{residual-stream}{residual stream}, and
the per-layer prediction tends to converge smoothly toward the
final-layer output. Concretely, for a pre-\glsterm{layer-normalization}{LayerNorm} \glsterm{transformer}{transformer} the
residual update at layer $\ell$ is
\begin{equation}
  \mathbf{h}_{\ell+1} \;=\; \mathbf{h}_\ell + F_\ell(\mathbf{h}_\ell),
  \label{eq:residual-update}
\end{equation}
where $F_\ell$ is the residual output of layer $\ell$ (attention
followed by MLP). Splitting the model $\mathcal{M}$ into the layers up
to $\ell$ and the layers after, the function from a \glsterm{hidden-state}{hidden state}
$\mathbf{h}_\ell$ to the final-layer \glsterm{logits}{logits} is
\citep[\S 2, Eq.~2]{belrose2023-tuned-lens}\footnote{KB excerpt:
\texttt{kb/excerpts/belrose2023-tuned-lens.md\#sec-2-decomposition}.}
\begin{equation}
  \mathcal{M}_{>\ell}(\mathbf{h}_\ell)
  \;=\;
  \mathrm{LayerNorm}\!\left[\mathbf{h}_\ell
  \;+\;\underbrace{\sum_{\ell'=\ell}^{L} F_{\ell'}(\mathbf{h}_{\ell'})}_{\text{residual update}}\right]\!W_U,
  \label{eq:residual-decomp}
\end{equation}
with $W_U \in \mathbb{R}^{D \times V}$ the unembedding matrix (often
weight-tied with the input embedding) and $\mathrm{LayerNorm}$ the
model's own final pre-unembed normalization. Tensor shapes for the
rest of this section: \glsterm{residual-stream}{residual stream} $\mathbf{h}_\ell \in \mathbb{R}^D$
at one $(B, S)$ position, with $D$ the model dim, $V = |\mathcal{V}|$
the \glsterm{vocabulary}{vocabulary}, and $L$ the layer count from Paper~1's notation. A
batched \glsterm{residual-stream}{residual stream} is $\BSH$;
the lens output is $\bsv$.

The \emph{prediction trajectory}
$\{\mathrm{Lens}_\ell(\mathbf{h}_\ell)\}_{\ell=0}^L$ across layers is
the diagnostic object: a single lens reading is a snapshot, but the
evolution across layers fingerprints the model's iterative computation.
\citet[\S 1]{belrose2023-tuned-lens} report that for clean
in-distribution prompts the trajectory converges smoothly to the final
prediction many layers before $L$, and that prompt-injected or
otherwise anomalous inputs deviate from the smooth trajectory in
detectable ways — the basis for a near-perfect prompt-injection
detector \citep[\S 5.3]{belrose2023-tuned-lens}\footnote{KB excerpt:
\texttt{kb/excerpts/belrose2023-tuned-lens.md\#sec-5-3}.}.

\subsection{The logit lens: no-training baseline}
\label{sec:logit-lens}

The \glsterm{logit-lens}{logit lens} drops the post-$\ell$ residual updates from
Eq.~\eqref{eq:residual-decomp} entirely — it sets
$\sum_{\ell'=\ell}^{L} F_{\ell'} = 0$ and reads
\citep[\S 2, Eq.~3]{belrose2023-tuned-lens}\footnote{KB excerpt:
\texttt{kb/excerpts/belrose2023-tuned-lens.md\#sec-2-logit-lens}.}:
\begin{equation}
  \boxed{\;\mathrm{LogitLens}(\mathbf{h}_\ell) \;:=\; \mathrm{LayerNorm}[\mathbf{h}_\ell]\,W_U\;}
  \label{eq:logit-lens}
\end{equation}
The result is a logit vector in $\mathbb{R}^V$; \glsterm{softmax}{softmax} gives the
layer-$\ell$ ``as-if-final'' distribution and $\arg\max$ gives the
top-$1$ token. Per-(layer, token) cost is one \glsterm{layer-normalization}{LayerNorm} plus one
$D \times V$ matmul — exactly the same compute as the model's actual
unembed step at the top of the stack. The technique has no training
\glsterm{parameters}{parameters} and no hyperparameters beyond the choice of normalization;
it is the cheapest interpretability readout of any kind.

\begin{forumsignal}
The \glsterm{logit-lens}{logit lens} was introduced by nostalgebraist on LessWrong in August
2020 \citep{nostalgebraist2020-logit-lens} and propagated for nearly
three years before any peer-reviewed treatment landed
\citep[\S 2]{belrose2023-tuned-lens}. We treat the LessWrong post as the
canonical source of the technique, but flag it explicitly as a
forum-tier (Tier C) artifact under the source-tier convention of
Paper~1: the technique is widely reused
\citep[Halawi 2023, Geva 2022, Dar 2022 referenced in
][\S 2]{belrose2023-tuned-lens}, but the original write-up is a
LessWrong essay, not a journal paper. Hard technical claims about lens
behavior in this section back to the Belrose et al.\ peer-reviewed
treatment, not to the original post.
\end{forumsignal}

\subsubsection*{Why the logit lens fails on some models}

\citet[\S 2]{belrose2023-tuned-lens} document three failure modes of
the raw \glsterm{logit-lens}{logit lens}\footnote{KB excerpt:
\texttt{kb/excerpts/belrose2023-tuned-lens.md\#sec-2-failures}.}.
\textbf{Bias.} The \glsterm{logit-lens}{logit lens} systematically puts more mass on certain
\glsterm{vocabulary}{vocabulary} items than the final layer does. On Pythia 12B, the bias
measured as KL between logit-lens marginal and final-layer marginal is
``around 4 to 5 bits for most layers,'' against $0.0068$ bits between
Pythia 160M's and Pythia 12B's final-layer
distributions~\citep[\S 2]{belrose2023-tuned-lens}. A rational agent's
beliefs should not update predictably over time, so the prediction
trajectory under a biased lens cannot be read as Bayesian belief
updating. \textbf{Representational drift.} Hidden states at different
layers occupy slightly different covariance subspaces; the final-layer
covariance often shifts sharply relative to earlier layers
\citep[\S 3]{belrose2023-tuned-lens}, so a fixed final-layer unembedding
applied to an earlier-layer state may be misapplied to the wrong basis.
\textbf{Cross-model unreliability.} The \glsterm{logit-lens}{logit lens} works ``reasonably
well'' on GPT-2 (\citealp{radford2019-gpt2}) but fails on BLOOM and
GPT-Neo: its top-$1$ predictions in the bottom 21 layers of GPT-Neo-2.7B
are nonsense (e.g., ``\texttt{@G:}'') rather than coherent intermediate
tokens \citep[\S 2 + Fig.~1]{belrose2023-tuned-lens}. For BLOOM and OPT
125M, the top-$1$ logit-lens prediction is the input token rather than
any plausible continuation in more than half the
layers~\citep[\S 2]{belrose2023-tuned-lens}.

These three failure modes set up the \glsterm{tuned-lens}{tuned lens} as the targeted
correction.

\subsection{The tuned lens: a per-layer affine translator}
\label{sec:tuned-lens}

\citet[\S 3]{belrose2023-tuned-lens} insert a learned affine
\emph{translator} $(A_\ell, \mathbf{b}_\ell)$ at each layer that maps
$\mathbf{h}_\ell$ into the basis the unembedding expects:
\begin{equation}
  \boxed{\;\mathrm{TunedLens}_\ell(\mathbf{h}_\ell)
  \;:=\; \mathrm{LogitLens}(A_\ell \mathbf{h}_\ell + \mathbf{b}_\ell)\;}
  \label{eq:tuned-lens}
\end{equation}
\citep[\S 3, Eq.~8]{belrose2023-tuned-lens}\footnote{KB excerpt:
\texttt{kb/excerpts/belrose2023-tuned-lens.md\#sec-3-tuned-lens}.},
with $A_\ell \in \mathbb{R}^{D \times D}$ and
$\mathbf{b}_\ell \in \mathbb{R}^D$. Translators are initialized to the
identity, so a freshly initialized \glsterm{tuned-lens}{tuned lens} is identical to the logit
lens; training pulls the translator off the identity to absorb basis
drift. Crucially, the unembedding $W_U$ is \emph{frozen} — there is
exactly one $W_U$ per model, shared across all $L$ tuned lenses
\citep[\S 3]{belrose2023-tuned-lens}.

The training objective is a per-layer distillation against the model's
own final-layer prediction:
\begin{equation}
  \mathcal{L}(\ell)
  \;=\;
  \mathbb{E}_\mathbf{x}\!\left[D_{\mathrm{KL}}\!\left(\mathcal{M}_{>\ell}(\mathbf{h}_\ell)
  \,\big\|\,\mathrm{TunedLens}_\ell(\mathbf{h}_\ell)\right)\right]
  \label{eq:tuned-lens-loss}
\end{equation}
\citep[\S 3, Eq.~9]{belrose2023-tuned-lens}\footnote{KB excerpt:
\texttt{kb/excerpts/belrose2023-tuned-lens.md\#sec-3-loss}.}.
Using $\mathcal{M}_{>\ell}$ as the soft target rather than ground-truth
tokens is what makes Eq.~\eqref{eq:tuned-lens-loss} a distillation loss:
the translator is incentivized to recover what the model itself already
knows, never to learn extra information beyond it. This is the design
choice that distinguishes tuned-lens translators from
\citet{alain-bengio-2017}-style early-exit probes, which retrain a full
unembedding per layer and risk learning structure not present in the
underlying model — the classic ``\glsterm{probe-2}{probe}-as-\glsterm{feature}{feature}-extractor'' confound
flagged by \citet{hewitt2019-control-tasks} and engaged in
\S\ref{sec:probing}.

\paragraph{Cost.} Tuned-lens \glsterm{parameters}{parameters} total $L \cdot (D^2 + D)$ — a
fraction of the model's own size, since $D \times D$ is roughly the
size of one attention block's $W^Q$ projection. For Pythia 12B with
$D = 5120$ and $L = 36$, this works out to roughly $944$M extra
\glsterm{parameters}{parameters} \citep[\S 3.2]{belrose2023-tuned-lens}\footnote{KB excerpt:
\texttt{kb/excerpts/belrose2023-tuned-lens.md\#sec-3-tuned-lens}.} ---
considerably smaller than per-layer unembeddings of size
$V \times D$ would be, since $V \in [50\text{K}, 250\text{K}]$ across
the models Belrose et al.\ study. The cost-per-(layer, token) at
inference is one extra $D \times D$ matmul on top of the logit-lens
\glsterm{layer-normalization}{LayerNorm} + $W_U$ chain.

\paragraph{Why a single affine map is enough.} The \glsterm{tuned-lens}{tuned lens} fixes
all three logit-lens failures with one affine map per layer:
$\mathbf{b}_\ell$ absorbs the bias term, $A_\ell$ absorbs the
representational drift (basis shift), and the joint training across all
layers adapts to the cross-model pathologies that defeated the raw
\glsterm{logit-lens}{logit lens}. \citet[Fig.~1]{belrose2023-tuned-lens} report that the
\glsterm{tuned-lens}{tuned lens} recovers coherent intermediate predictions throughout the
GPT-Neo-2.7B stack where the \glsterm{logit-lens}{logit lens} produced nonsense, and
generalizes across Pythia, GPT-Neo, GPT-NeoX-20B, OPT, BLOOM, and
GPT-2. The fact that a \emph{linear} correction suffices is itself
informative: representational drift in residual-stream architectures is
approximately linear, not learned-\glsterm{feature}{feature}-dependent.

\begin{intuition}
The lens framework exploits one architectural fact: the \glsterm{residual-stream}{residual stream}
is an additive bus. The final-layer \glsterm{hidden-state}{hidden state} is
$\mathbf{h}_L = \mathbf{h}_\ell + \sum_{\ell'=\ell}^{L-1} F_{\ell'}(\mathbf{h}_{\ell'})$
by recursing Eq.~\eqref{eq:residual-update}. If the residual sum is
small relative to $\mathbf{h}_\ell$, the unembedding of $\mathbf{h}_\ell$
already approximates the unembedding of $\mathbf{h}_L$ — the \glsterm{logit-lens}{logit lens}
is a zero-order Taylor expansion that the residual is unchanged. The
tuned-lens translator
$A_\ell \mathbf{h}_\ell + \mathbf{b}_\ell$ in
Eq.~\eqref{eq:tuned-lens} extends this to a first-order correction
that compensates for whatever systematic drift the actual residual
contributes at layer $\ell$. Returning to canonical form: the tuned
lens is the \glsterm{logit-lens}{logit lens} with a per-layer linear basis correction; the
underlying readout is still $\mathrm{LayerNorm}[\,\cdot\,]W_U$.
\end{intuition}

\subsection{DoLa: lens techniques as a decoding strategy}
\label{sec:dola}

The lens framework is not only diagnostic. \citet{chuang2023-dola}
turn the layerwise readout into an inference-time decoding strategy by
contrasting \glsterm{logits}{logits} across layers. \emph{Decoding by Contrasting Layers}
(\glsterm{dola}{DoLa}) defines, at each generation step,
\begin{equation}
  \mathrm{logits}_{\mathrm{DoLa}}(t)
  \;=\;
  \mathrm{logit}^{(L)}(t) \;-\; \mathrm{logit}^{(\ell^\star)}(t),
  \label{eq:dola}
\end{equation}
where $\ell^\star$ is selected per token as the layer whose logit
distribution is maximally divergent (in Jensen-Shannon divergence) from
the final-layer logit distribution \citep{chuang2023-dola}. Tokens
whose probability rises sharply between layer $\ell^\star$ and layer
$L$ are exactly the tokens the later layers specifically committed to;
reweighting toward those tokens preserves the late-stage factual
signal while suppressing tokens that were already high in early layers
(plausibly hallucination-shaped). \glsterm{dola}{DoLa} requires no fine-tuning, adds
only a modest decoding-time overhead, and reports improved TruthfulQA
scores on LLaMA-7B/13B/33B \citep{chuang2023-dola}. The conceptual
point: layerwise readout carries decision-relevant information that
can be fed back into generation, not just inspected.

\begin{analogy}
\glsterm{dola}{DoLa} is to the \glsterm{logit-lens}{logit lens} what a numerical \emph{difference} is to a
function evaluation. The \glsterm{logit-lens}{logit lens} reads
$\mathrm{LogitLens}(\mathbf{h}_\ell)$ as the model's belief at layer
$\ell$; \glsterm{dola}{DoLa} reads
$\mathrm{LogitLens}(\mathbf{h}_L) - \mathrm{LogitLens}(\mathbf{h}_{\ell^\star})$
as the \emph{change} in belief that the later layers contributed. The
analogy returns to canonical form: Eq.~\eqref{eq:dola} is exactly the
difference of two logit-lens evaluations
(Eq.~\eqref{eq:logit-lens}) on the same token, so the underlying
object is still the layerwise unembed — read differentially across
depth instead of pointwise at one layer.
\end{analogy}

\subsection{Diagnostic: the prediction trajectory}
\label{sec:lens-trajectory}

The empirical signature of a working lens is the smoothness of the
prediction trajectory $\{\mathrm{Lens}_\ell(\mathbf{h}_\ell)\}_{\ell=0}^L$.
\citet[\S 1]{belrose2023-tuned-lens} document three regimes. On
\emph{clean in-distribution prompts}, the tuned-lens trajectory
converges smoothly to the final prediction, with each layer typically
reducing \glsterm{perplexity}{perplexity}. On \emph{prompt-injected} or otherwise anomalous
inputs, the trajectory deviates from the smooth manifold of normal
trajectories detectably enough that off-the-shelf isolation-forest /
local-outlier-factor classifiers achieve near-perfect prompt-injection
detection accuracy \citep[\S 5.3]{belrose2023-tuned-lens}\footnote{KB
excerpt: \texttt{kb/excerpts/belrose2023-tuned-lens.md\#sec-5-3}.}. On
\emph{hard examples} — data points the model required many training
steps to learn — the trajectory's commit-layer extends deeper into the
stack: harder examples are classified later
\citep[\S 5.4]{belrose2023-tuned-lens}. The trajectory is therefore
not just a richer readout than the layer-pointwise prediction; it is
also a fingerprint that distinguishes input regimes the model treats
differently.

\begin{contradiction}
\textbf{Layer-of-prediction vs.\ layer-of-decision.} The trajectory
view assumes the lens reading is monotonically improving in $\ell$:
later layers reflect the model's better-considered prediction.
\citet{halawi2023-overthinking}, however, find that on certain tasks
predictions extracted from \emph{earlier} layers are \emph{more}
robust to incorrect demonstrations than the final-layer prediction.
\citet[\S 5.2]{belrose2023-tuned-lens} reproduce this on three models
(BLOOM 560M, GPT-Neo 1.3B, GPT-Neo 2.7B), but flag that without
peeking at ground-truth labels the lens user has no a-priori recipe
for choosing the right layer\footnote{KB excerpt:
\texttt{kb/excerpts/belrose2023-tuned-lens.md\#sec-5}.}. So
``commit later $=$ commit better'' is sometimes false, and the field
lacks a principled criterion for layer selection. The strongest
form of the claim is that lens techniques diagnose what \emph{would}
happen at $\ell$, not what the model \emph{should} commit to; closing
that gap requires the causal interventions of \S\ref{sec:probing}
and the patching machinery of \S\ref{sec:patching}.
\end{contradiction}

\subsection{Frontier coverage and forward link}
\label{sec:lens-forward}

The lens framework as published covers the open-source generation
through 2023: GPT-2, Pythia (70M–12B), GPT-Neo, GPT-NeoX-20B, OPT,
BLOOM, and (via lens transfer) Vicuna-13B
\citep[\S 3, \S 5]{belrose2023-tuned-lens}. The 2025–2026 frontier ---
Llama 3 / 4, Qwen 2.5 / 3, \glsterm{precision-pinning}{DeepSeek-V3}, Gemma 3 --- is largely
uncovered by published tuned lenses as of writing. The
\texttt{logitlens4llms-2025} extension targets Qwen-2.5 and Llama-3.1;
broader open-lens releases for the rest of the frontier-2026 lineup
are an outstanding need. Whether the lens framework extends cleanly
to architectures with substantially modified attention or \glsterm{ffn}{FFN} substrate
--- \glsterm{precision-pinning}{DeepSeek-V3}'s \glsterm{multi-head-latent-attention}{MLA}-compressed-KV plus fine-grained-MoE, \glsterm{nsa}{NSA}-sparse
attention, sliding-window layers --- is partially empirical and not
yet documented at frontier scale. The \glsterm{residual-stream}{residual stream} is still
well-defined in all of these, but the layer structure is more
irregular, and the assumption that a single affine translator suffices
to absorb basis drift may need re-examination on per-architecture
basis.

The lens is correlational by construction: it tells you what the model
\emph{would} predict if it committed at $\ell$, not whether the model
\emph{uses} the layer-$\ell$ information that supports that prediction.
\citet[\S 4]{belrose2023-tuned-lens} partially close this gap with
\emph{\glsterm{causal-basis-extraction}{Causal Basis Extraction}}, which finds a tuned-lens-derived
orthonormal basis ordered by mean-ablation influence and validates that
the directions important to the lens are also important to the
underlying model (Spearman $\rho = 0.89$ on Pythia 410M, layer
$18$)\footnote{KB excerpt:
\texttt{kb/excerpts/belrose2023-tuned-lens.md\#sec-4-cbe}.}. The next
two sections take up this correlational/causal split directly: §%
\ref{sec:probing} introduces probes as the supervised correlational
analog of lenses --- a trained classifier on frozen activations rather
than the model's own unembedding --- and §\ref{sec:correlation-causation}
shows how the field couples probes with activation interventions to
upgrade correlational readouts into causal claims. The lens results in
this section reappear in §\ref{sec:cross-method}: the tuned-lens
prediction trajectory on IOI prompts is one of the three independent
methodologies that converge on the GPT-2-small IOI-\glsterm{circuit}{circuit} picture, and
the lens cross-model unreliability documented in
\S\ref{sec:logit-lens} is one of the field's open contradictions
catalogued there.
